\documentclass[11pt]{article}
\usepackage[top=0.85in,bottom=0.85in,left=0.9in,right=0.9in]{geometry}
\usepackage{amsmath,amssymb}
\usepackage{booktabs,multirow}
\usepackage{hyperref}
\usepackage{microtype}
\usepackage{enumitem}
\setlist{nosep,leftmargin=*,topsep=2pt}
\setlength{\parskip}{3pt}
\setlength{\parindent}{1em}

\title{\textbf{Re-implementing Cold Diffusion on CIFAR-10}\\[2pt]
       \large CS 4782 Final Project Report}
\author{Johnson Wang (jw2693) \quad Eric Weng (ew522) \\ Cornell University}
\date{May 12, 2026}

\begin{document}
\maketitle
\vspace{-1.2em}

\section{Introduction}

Standard diffusion models iteratively add Gaussian noise and train a neural network to
reverse it. \emph{Cold Diffusion: Inverting Arbitrary Image Transforms Without
Noise}~\cite{bansal2022} challenges this design: Gaussian noise is not a necessary
ingredient. Any smoothly-parameterized, deterministic image degradation can replace it,
and the resulting model retains the ability to reconstruct and synthesize images. The key
contribution is \textbf{Algorithm~2} (§3.3 of the paper), an anchored reverse sampler
that at step~$s$ computes $x_{s-1} = x_s - D(\hat{x}_0, s) + D(\hat{x}_0, s{-}1)$
where $\hat{x}_0 = R_\theta(x_s, s)$, stabilizing the trajectory for any choice of
deterministic forward process~$D$.

We re-implement Cold Diffusion from scratch in PyTorch on CIFAR-10, reproducing all
three degradations from the paper (inpainting, Gaussian blur, super-resolution) using
the published training recipe. We additionally extend the paper with
\textbf{Algorithm~3}, an EMA-smoothed variant of Algorithm~2, and validate it on both
CIFAR-10 and MNIST.

\section{Chosen Result}

We target two results: \textbf{Tables~1--2} (FID / SSIM / RMSE for degraded, direct, and
sampled outputs on CIFAR-10 inpainting and blur) and \textbf{Figure~2} (the qualitative
sampler comparison showing Algorithm~1 failing and Algorithm~2 succeeding). These jointly
establish the paper's central claim: the anchored update in Algorithm~2 is essential for
smooth degradations. For inpainting the paper reports a sampled FID of 8.92 vs.\ a degraded
FID of 40.83, giving a clear quantitative target. We also provide Algorithm~1 metrics
(not in the paper) to quantify the failure mode.

\section{Methodology}

\paragraph{U-Net.}
3-level U-Net with base channels 64, multipliers $(1,2,2)$, two residual blocks per level
(GroupNorm + SiLU activations), and self-attention at the $8{\times}8$ bottleneck.
Sinusoidal time embeddings projected through a 2-layer MLP are injected into each residual
block. Total: ${\approx}12$M parameters.

\paragraph{Degradations (Appendix A).}
\begin{itemize}
  \item \emph{Inpainting}: cumulative Gaussian mask
    $m_{\beta_t}(i,j)=1-\exp(-d^2/2\beta_t^2)$, $\beta_1{=}1$, $\beta_{t+1}{=}\beta_t+0.1$,
    $T{=}50$, center randomized per image on CIFAR-10.
  \item \emph{Blur}: 11${\times}$11 Gaussian kernel applied recursively,
    $\sigma_t = 0.01t+0.35$, $T{=}40$.
  \item \emph{Super-resolution}: $2{\times}$ avg-pool + nearest upsample, $T{=}3$.
\end{itemize}

\paragraph{Training.}
L1 loss $\mathbb{E}_{t,x_0}[\|R_\theta(D(x_0,t),t)-x_0\|_1]$; Adam at $2{\times}10^{-5}$;
batch~64 with 2-step gradient accumulation (effective batch~128); EMA decay~0.995 updated
every 10 steps. Inpainting: 60k steps; blur: 30k steps (paper uses 700k); SR: 100k steps.
Hardware: single NVIDIA T4 GPU on Google Colab.

\paragraph{Algorithm~3 (our contribution).}
Algorithm~2 discards each per-step $\hat{x}_0$ prediction immediately. We maintain instead
a running EMA: $\hat{x}_0^{(s)} \leftarrow \alpha\,\hat{x}_0^{(s+1)} + (1{-}\alpha)\,R(x_s,s)$,
then use the smoothed $\hat{x}_0^{(s)}$ in the anchored update. Setting $\alpha{=}0$ recovers
Algorithm~2. We sweep $\alpha \in \{0.1, 0.3, 0.5, 0.7, 1.0\}$ on the CIFAR-10 test set.

\paragraph{Key implementation note.}
For randomized inpainting, every $D(\cdot)$ call within one trajectory must share the same
per-image mask center; otherwise the subtraction $D(\hat{x}_0,s) - D(\hat{x}_0,s{-}1)$ uses
mismatched masks and produces artifacts. We implement explicit \emph{state-passing}: the center
is drawn once per trajectory and locked for all subsequent $D$ calls. This requirement is
implicit in the paper and was the primary debugging challenge.

\section{Results \& Analysis}

Table~\ref{tab:metrics} shows FID / SSIM / RMSE for inpainting (60k steps) and blur (30k
steps). All Algorithm~1 numbers are ours; the paper does not report them.

\begin{table}[h!]
\centering
\caption{Reconstruction metrics on CIFAR-10. Lower FID/RMSE and higher SSIM are better.
         Paper values from Tables~1--2~\cite{bansal2022}. Alg~1 numbers ours only.}
\label{tab:metrics}
\vspace{3pt}
\small
\setlength{\tabcolsep}{5pt}
\begin{tabular}{ll rr rr rr rr}
\toprule
 & & \multicolumn{2}{c}{Degraded} & \multicolumn{2}{c}{Direct} &
   \multicolumn{2}{c}{Alg 1} & \multicolumn{2}{c}{Alg 2 Sampled} \\
\cmidrule(lr){3-4}\cmidrule(lr){5-6}\cmidrule(lr){7-8}\cmidrule(lr){9-10}
Task & Metric & Paper & Ours & Paper & Ours & Paper & Ours & Paper & Ours \\
\midrule
\multirow{3}{*}{\shortstack[l]{Inpainting\\(T=50)}}
  & FID$\downarrow$  & 40.83 & 70.73 &  9.71 & 15.21 & — & 123.17 &  8.92 & 12.16 \\
  & SSIM$\uparrow$   & 0.615 & 0.579 & 0.869 & 0.877 & — &  0.440 & 0.859 &  0.807 \\
  & RMSE$\downarrow$ & 0.143 & 0.302 & 0.063 & 0.071 & — &  0.248 & 0.068 &  0.081 \\
\midrule
\multirow{3}{*}{\shortstack[l]{Blur\\(T=40)}}
  & FID$\downarrow$  & 238.26 & 252.71 & 83.69 & 99.18 & — & 105.98 & 80.08 & 96.59 \\
  & SSIM$\uparrow$   &  0.315 &  0.410 &  0.875 & 0.789 & — &  0.621 &  0.873 & 0.794 \\
  & RMSE$\downarrow$ &  0.136 &  0.125 &  0.071 & 0.062 & — &  0.112 &  0.075 & 0.063 \\
\bottomrule
\end{tabular}
\end{table}

\noindent\textbf{Algorithm 1 vs.\ 2.}
The most striking result is how completely Algorithm~1 fails: inpainting FID jumps to
123.17 (vs.\ Algorithm~2's 12.16) and SSIM collapses to 0.440. Each naive step re-degrades
a fresh prediction using a new random mask center, compounding errors and erasing spatial
coherence. Algorithm~2's anchored update eliminates this entirely.

\noindent\textbf{Gap from paper.}
Our inpainting results (FID 12.16 vs.\ paper's 8.92) are close; the gap is consistent with
a slightly smaller U-Net---the paper does not fully specify its architecture. For blur the
gap is larger (FID 96.59 vs.\ 80.08) because we trained only 30k of the paper's 700k steps;
the Alg~1 vs.\ Alg~2 trend is still clearly visible.

\noindent\textbf{Algorithm 3.}
Table~\ref{tab:alg3} compares Alg~2 and Alg~3 at the best $\alpha{=}0.3$
(FID sweep: 12.44, \textbf{11.89}, 12.15, 12.45, 13.62 for $\alpha \in \{0.1, 0.3, 0.5, 0.7, 1.0\}$).

\begin{table}[h!]
\centering
\caption{Algorithm~2 vs.\ Algorithm~3 ($\alpha{=}0.3$) on CIFAR-10.}
\label{tab:alg3}
\vspace{3pt}
\small
\begin{tabular}{llrrr}
\toprule
Task & Method & FID$\downarrow$ & SSIM$\uparrow$ & RMSE$\downarrow$ \\
\midrule
\multirow{2}{*}{Inpainting (T=50)} & Alg 2 & 12.16 & 0.807 & 0.081 \\
                                    & Alg 3 & \textbf{11.89} & \textbf{0.872} & \textbf{0.080} \\
\midrule
\multirow{2}{*}{Blur (T=40)}       & Alg 2 & 96.59 & 0.794 & 0.063 \\
                                    & Alg 3 & 97.37 & 0.794 & 0.062 \\
\bottomrule
\end{tabular}
\end{table}

\noindent On inpainting (nonlinear $D$) Algorithm~3 improves FID by 0.27 and SSIM by 0.065.
On blur (linear $D$) Algorithm~3 essentially ties. This asymmetry precisely confirms the
hypothesis: the EMA helps when per-step $D$ calls are only approximately consistent
(random mask centers make inpainting weakly nonlinear in the trajectory), but adds no signal
when $D$ is exactly linear and Algorithm~2 is already optimal (§3.3).

\section{Reflections}

The biggest practical challenge was the state-locking requirement for randomized inpainting.
Without fixing the mask center for the entire trajectory, Algorithm~2's subtraction
$D(\hat{x}_0, s) - D(\hat{x}_0, s{-}1)$ uses inconsistent masks and the output is visually
indistinguishable from Algorithm~1's failures---a subtle bug that took significant debugging
to isolate. We solved it with an explicit \texttt{sample\_state()} abstraction that draws
the stochastic part of $D$ once per trajectory and passes it to all subsequent calls.

The broader lesson is that the paper's theoretical analysis in §3.3 (Algorithm~2 is exact
for linear $D$) is tight enough to guide empirical design. Algorithm~3's clean
linear/nonlinear split confirms that EMA smoothing is useful specifically where the theory
predicts imperfection. Given more compute we would: (1) train blur/SR to the paper's full
step budgets; (2) anneal $\alpha$ during sampling rather than holding it fixed; (3) test on
CelebA (256${\times}$256) where inpainting's spatial nonlinearity is stronger.

\begin{thebibliography}{9}
\bibitem{bansal2022}
  Bansal, A., et al.
  \emph{Cold Diffusion: Inverting Arbitrary Image Transforms Without Noise.}
  NeurIPS 2022. \href{https://arxiv.org/abs/2208.09392}{arXiv:2208.09392}.
\bibitem{ho2020}
  Ho, J., Jain, A., Abbeel, P.
  \emph{Denoising Diffusion Probabilistic Models.}
  NeurIPS 2020.
\bibitem{song2021}
  Song, J., Meng, C., Ermon, S.
  \emph{Denoising Diffusion Implicit Models.}
  ICLR 2021.
\bibitem{krizhevsky2009}
  Krizhevsky, A.
  \emph{Learning Multiple Layers of Features from Tiny Images.}
  Technical Report, 2009.
\bibitem{pytorch}
  Paszke, A., et al.\ PyTorch. \texttt{torchvision}, \texttt{torchmetrics}.
\end{thebibliography}

\end{document}
